\documentclass[11pt,a4paper]{article}

\usepackage[spanish,es-nodecimaldot]{babel}
\usepackage[utf8]{inputenc}
\usepackage[T1]{fontenc}
\usepackage{lmodern}
\usepackage{geometry}
\geometry{top=2.3cm, bottom=2.3cm, left=2.5cm, right=2.5cm}

\usepackage{xcolor}
\definecolor{arqmedDark}{HTML}{0F172A}
\definecolor{arqmedAmber}{HTML}{D97706}
\definecolor{arqmedSlate}{HTML}{334155}
\definecolor{arqmedMuted}{HTML}{64748B}
\definecolor{arqmedLight}{HTML}{F8FAFC}
\definecolor{arqmedBorder}{HTML}{CBD5E1}

\usepackage{hyperref}
\hypersetup{
    colorlinks=true,
    linkcolor=arqmedDark,
    urlcolor=arqmedAmber,
    citecolor=arqmedSlate,
    pdftitle={Manual de Usuario - ARQMED Studio App Tablet},
    pdfauthor={ARQMED Team}
}

\usepackage{tcolorbox}
\tcbuselibrary{skins,breakable}

\newtcolorbox{pasobox}[1][Paso a paso]{
    enhanced,
    breakable,
    colback=white,
    colframe=arqmedDark,
    arc=2.5mm,
    title={\textbf{#1}},
    fonttitle=\bfseries\color{white},
    coltitle=white,
    colbacktitle=arqmedDark,
    left=4mm, right=4mm, top=3mm, bottom=3mm,
    boxrule=0.8pt
}

\newtcolorbox{aviso}[1][Atención en obra]{
    enhanced,
    breakable,
    colback=arqmedLight,
    colframe=arqmedAmber,
    arc=2.5mm,
    title={\textbf{#1}},
    fonttitle=\bfseries\color{white},
    coltitle=white,
    colbacktitle=arqmedAmber,
    left=4mm, right=4mm, top=3mm, bottom=3mm,
    boxrule=0.8pt
}

\usepackage{booktabs}
\usepackage{tabularx}
\usepackage{enumitem}
\usepackage{titlesec}

\titleformat{\section}
  {\normalfont\Large\bfseries\color{arqmedDark}}{\thesection}{1em}{}[\titlerule]
\titleformat{\subsection}
  {\normalfont\large\bfseries\color{arqmedSlate}}{\thesubsection}{1em}{}

\begin{document}

\begin{titlepage}
    \centering
    \vspace*{1.5cm}
    
    {\Huge \textbf{\color{arqmedDark}ARQMED Studio}}\\[0.4cm]
    {\LARGE \textbf{\color{arqmedAmber}Guía de Uso de la Aplicación en Tablet}}\\[0.4cm]
    {\large Manual práctico para consultar, anotar mediciones y editar presupuestos .BC3 en movilidad}\\
    
    \vspace{2cm}
    
    \begin{tcolorbox}[colback=arqmedLight,colframe=arqmedDark,arc=3mm,width=0.9\textwidth,center]
        \centering
        \vspace{0.3cm}
        \textbf{Dispositivo objetivo:} Tabletas Android y pantallas táctiles\\
        \textbf{Instalador:} \texttt{arqmed-studio-v1.0.0.apk}\\
        \textbf{Funcionamiento:} 100\% fuera de línea (sin cobertura ni conexión a internet)\\
        \vspace{0.3cm}
    \end{tcolorbox}
    
    \vfill
    
    {\color{arqmedMuted}
    Instrucciones directas: instalación, apertura de archivos propios,\\
    gestión de la biblioteca local, toma de datos en obra y guardado final.\\[0.5cm]
    \today
    }
\end{titlepage}

\tableofcontents
\newpage

\section{Instalación en la tableta}

La aplicación se distribuye como un paquete instalable de Android (.apk) que no depende de Google Play ni requiere cuentas de usuario.

\begin{pasobox}[Cómo instalar la aplicación]
\begin{enumerate}[leftmargin=*]
    \item Descarga el archivo \texttt{arqmed-studio-v1.0.0.apk} desde la sección de Releases del proyecto en GitHub.
    \item Pásalo a la tableta mediante cable USB, tarjeta microSD, Google Drive o Telegram.
    \item Abre el gestor de archivos de tu tableta (la aplicación «Archivos» o «Mis Archivos») y pulsa sobre el archivo descargado.
    \item Si Android te advierte sobre la instalación de aplicaciones desconocidas, pulsa en \textbf{Ajustes} y activa la casilla para autorizar la instalación desde esa fuente.
    \item Pulsa en \textbf{Instalar} y, al terminar, en \textbf{Abrir}.
\end{enumerate}
\end{pasobox}

En el menú de aplicaciones de tu tableta aparecerá el icono con el isotipo de ARQMED (con sus colores corporativos en blanco, pizarra y naranja técnico sobre fondo blanco).

\section{Primer contacto con la pantalla}

Al abrir la aplicación verás una pantalla de inicio oscura con el logotipo centrado que dura un segundo mientras se prepara la memoria local. Inmediatamente después se carga la interfaz en modo claro (blanco), pensada para trabajar cómodamente incluso con luz natural en obra.

La pantalla de la tableta se divide en tres niveles:

\subsection{1. La barra superior}
Diseñada para no ocupar espacio útil en pantallas compactas:
\begin{itemize}[leftmargin=*]
    \item \textbf{A la izquierda}: el logotipo de ARQMED y el nombre del archivo presupuestario que tienes abierto en ese momento. Al lado verás una pequeña pastilla verde que indica «✓ Local» o la hora del último cambio.
    \item \textbf{A la derecha}: un único botón destacado con el texto \textbf{📁 Mis .BC3}. Al pulsarlo se abre el gestor donde puedes ver todos los presupuestos que tengas guardados en la tableta o importar uno nuevo.
\end{itemize}

\subsection{2. La barra de menús}
Justo debajo se sitúan los menús desplegables adaptados a toques de dedo:
\begin{itemize}[leftmargin=*]
    \item \textbf{Archivo}: permite guardar el presupuesto en la tableta, importar archivos nuevos, duplicar obras o exportar copias a CSV (Excel).
    \item \textbf{Editar}: opciones para deshacer (\texttt{Ctrl+Z}), rehacer, editar la partida seleccionada o añadir conceptos nuevos.
    \item \textbf{Ver}: desde aquí cambias entre el \textbf{Editor árbol} (la vista habitual), la \textbf{Matriz multi-nivel} y el \textbf{Dashboard de obra} con gráficos analíticos. También puedes plegar o desplegar capítulos por niveles con un solo toque.
    \item \textbf{Preferencias}: permite alternar al tema oscuro si trabajas de noche o en interiores con poca luz.
\end{itemize}

\subsection{3. El área central y los paneles de apoyo}
La zona principal muestra el árbol completo del presupuesto (capítulos, subcapítulos y partidas). En la parte inferior se despliegan automáticamente los paneles de apoyo:
\begin{itemize}[leftmargin=*]
    \item \textbf{Descomposición}: muestra los recursos elementales de la partida (materiales, mano de obra, maquinaria).
    \item \textbf{Mediciones}: tabla con las líneas de cotas (unidades, longitud, anchura y altura).
    \item \textbf{Pliego / Comentario}: texto técnico completo de la partida (registro \texttt{\textasciitilde T}).
\end{itemize}

\section{Cómo trabajar con tus propios archivos .BC3}

Cuando abres la aplicación por primera vez verás cargado un presupuesto de demostración para que puedas explorar la herramienta sin esperas. Para cargar el presupuesto real de tu proyecto:

\subsection{Pasar el archivo a la tableta}
Copia tu archivo con extensión \texttt{.bc3} (exportado desde Presto, Arquímedes, CYPE o Menfis) en cualquier carpeta accesible de la tableta, por ejemplo en la carpeta \textbf{Descargas} o en \textbf{Documentos}.

\subsection{Importar el archivo a ARQMED Studio}
\begin{enumerate}[leftmargin=*]
    \item Pulsa en la barra superior en el botón \textbf{📁 Mis .BC3} (o entra en \textbf{Archivo $\rightarrow$ Importar .BC3...}).
    \item Se abrirá el selector oficial de archivos de Android.
    \item Busca tu archivo y pulsa sobre él.
    \item El motor leerá el presupuesto, convertirá automáticamente las codificaciones antiguas si procede y lo registrará en la biblioteca interna de la aplicación.
    \item La pantalla se actualizará de inmediato mostrando tu obra lista para trabajar.
\end{enumerate}

\begin{aviso}[¿Qué pasa con la confidencialidad de tu presupuesto?]
Tu archivo .BC3 se procesa de forma íntegra dentro de la tableta (en la base de datos interna del navegador). No se sube a internet, no requiere conexión WiFi ni datos móviles, y nadie ajeno al dispositivo puede acceder a tus mediciones o precios.
\end{aviso}

\section{Toma de mediciones y revisión en obra}

El uso principal de la tableta es revisar tajos en obra y comprobar mediciones sobre el terreno:

\subsection{Desplegar y buscar partidas}
\begin{itemize}[leftmargin=*]
    \item Toca sobre la flecha de cualquier capítulo para desplegar sus subcapítulos.
    \item Cada nivel tiene un distintivo visual propio (Capítulo general, Subcapítulo, Partida verde).
    \item Si quieres desplegar rápidamente todas las partidas de la obra, pulsa en \textbf{Ver $\rightarrow$ Selector de profundidad} y elige el nivel que quieras visualizar de un golpe.
\end{itemize}

\subsection{Anotar o corregir cotas de medición}
\begin{enumerate}[leftmargin=*]
    \item Pulsa sobre la partida que estés inspeccionando en obra.
    \item En la parte inferior de la pantalla, pulsa en la pestaña \textbf{Mediciones}.
    \item Verás las líneas de medición existentes con sus columnas de comentario, unidades, longitud, anchura y altura.
    \item Pulsa sobre cualquier celda numérica para que aparezca el teclado de la tableta y escribe la nueva cota.
    \item Si necesitas añadir una nueva medición de un tramo ejecutado hoy, pulsa el botón \textbf{+ Añadir línea} e introduce la zona y las medidas.
    \item Al cambiar cualquier cota, el subtotal de la línea y la medición acumulada de la partida se recalculan en tiempo real.
\end{enumerate}

\subsection{Modificar precios y controlar el impacto (OK / NO)}
Si negocias un precio nuevo de un material o cambias el rendimiento de una cuadrilla:
\begin{enumerate}[leftmargin=*]
    \item Pulsa en la partida y luego pulsa en \textbf{Editar concepto} en el menú o haz doble toque sobre el importe.
    \item Introduce el nuevo precio unitario.
    \item Si ese material interviene en más partidas de la obra, se abrirá la ventana de \textbf{Alcance de Impacto}.
    \item La ventana te mostrará un listado de todas las partidas donde participa ese material y cuánto dinero sube o baja cada una.
    \item Marca \textbf{OK} en las partidas donde quieras aplicar el cambio de precio y \textbf{NO} donde prefieras congelar el precio original.
    \item Pulsa en confirmar: el total de la partida y el importe global del presupuesto se actualizarán al instante hacia la raíz.
\end{enumerate}

\section{Guardar y llevarse el presupuesto modificado}

Mientras trabajas en la tableta, cada cambio que haces se va guardando automáticamente segundo a segundo en la memoria interna de la aplicación. No hay riesgo de perder datos si la tableta se apaga por batería o si cierras la app.

Cuando quieras extraer el archivo final para enviarlo a la oficina, a la constructora o abrirlo en el ordenador del despacho:

\begin{pasobox}[Cómo descargar el archivo .BC3 modificado]
\begin{enumerate}[leftmargin=*]
    \item Pulsa en la barra de menús en \textbf{Archivo}.
    \item Pulsa en la primera opción: \textbf{Guardar .BC3}.
    \item La aplicación empaqueta todos los cambios (mediciones nuevas, textos de pliego y precios) y genera un archivo con el nombre:
    \begin{center}
    \texttt{NOMBRE\_DE\_TU\_PRESUPUESTO\_EDITADO.bc3}
    \end{center}
    \item El archivo se descargará en la carpeta \textbf{Download} (Descargas) de tu tableta.
    \item Desde esa carpeta puedes enviarlo por correo electrónico, compartirlo por WhatsApp o copiarlo por USB a tu ordenador de sobremesa para abrirlo con Presto o Arquímedes sin ningún error de estructura.
\end{enumerate}
\end{pasobox}

\section{Cambiar entre varios presupuestos}

Si llevas el control de varias obras a la vez (por ejemplo, dos reformas y una vivienda unifamiliar):
\begin{itemize}[leftmargin=*]
    \item Pulsa el botón \textbf{📁 Mis .BC3} arriba a la derecha.
    \item Verás una tarjeta por cada presupuesto que hayas importado.
    \item Cada tarjeta muestra el título del proyecto, el importe total y la fecha de última modificación.
    \item Pulsa en \textbf{✓ Usar este .BC3} en la tarjeta de la obra que vayas a visitar para conmutar de una a otra en un instante.
    \item Si quieres hacer pruebas de costes sin alterar el presupuesto oficial, pulsa en \textbf{Duplicar}: se creará una copia gemela independiente para estudiar variantes.
\end{itemize}

\section{Consejos prácticos para uso en obra}

\begin{itemize}[leftmargin=*]
    \item \textbf{Modo avión y ahorro de batería}: en visitas largas a obra puedes poner la tableta en modo avión para ahorrar batería; ARQMED Studio funciona exactamente igual con o sin conexión.
    \item \textbf{Plegado de paneles}: si quieres ver muchas partidas a la vez en la pantalla vertical de la tableta, puedes arrastrar la barra divisoria inferior hacia abajo para dar más altura al árbol de capítulos.
    \item \textbf{Exportación rápida a Excel}: si el cliente o el encargado te pide una lista rápida de partidas en obra y no tiene un programa de presupuestos, entra en \textbf{Archivo $\rightarrow$ Exportar CSV (Excel)} y se generará una tabla que abre directamente en cualquier visor de hojas de cálculo de la tableta.
\end{itemize}

\end{document}
